\documentclass[10pt,a4paper]{amsart}
\usepackage[margin=19mm]{geometry}
\usepackage{amsmath,amssymb,mathtools}
\usepackage[scr=rsfs,cal=euler]{mathalfa}
\usepackage{xfrac}
\usepackage[unicode,hidelinks,bookmarks=false]{hyperref}
\newcommand{\ie}{i.e.\ }
\newcommand{\cf}{cf.\ }
\newcommand{\resp}{respectively\ }
\newcommand{\Subsection}{\subsection}
\providecommand{\nobreakdash}{\nobreak\hbox{-}}
\setlength{\emergencystretch}{2em}
\multlinegap0pt
\usepackage{amssymb}
\usepackage{mathtools}
\usepackage{xcolor}
\newtheorem{theorem}{Theorem}
\title{Error of discretization of Caputo fractional derivative in weighted spaces}
\author{\L{}ukasz P\l{}ociniczak$^*$, \and Hubert Woszczek$^{*, }$}
\date{Source: arXiv:2604.22470v1 (CC BY 4.0)}
\begin{document}
\maketitle

\noindent\emph{Source and licence.} Error of discretization of Caputo fractional derivative in weighted spaces. Version arXiv:2604.22470v1 (CC BY 4.0), distributed by arXiv under the Creative Commons Attribution 4.0 International licence, http://creativecommons.org/licenses/by/4.0/, as recorded in the arXiv licence metadata for this exact version and shown on the abstract page https://arxiv.org/abs/2604.22470v1. Full licence notice: \texttt{licenses/arxiv-2604.22470-CC-BY-4.0.txt}.\par\smallskip

\noindent\textit{Editorial scope.} Counted unit: the article's theorem th1 (source0.tex lines 125--134). The article's complete original proof of it is delivered with it (source0.tex lines 136--198, one \texttt{proof} environment of 63 lines). No mathematics is written, repaired or re-derived: the statement and the proof are the article's own lines, in the article's own notation, and no retained line is substituted or re-flowed.  No further source-local context is needed: every cross-reference of every retained line resolves inside the two delivered spans.  Nothing was omitted from any retained span, so the package carries no omission record.  No retained line contains a citation, so the wrapper carries no bibliography and no bibliography entry.  No retained line contains a figure environment, an image-inclusion command or drawing code, so no image is delivered and the package declares no asset dependency.  Editorial formatting only, in addition: an AMS \texttt{amsart} preamble that restates the article's own declarations and macros used by the retained lines, namely the environments \texttt{theorem} on separate counters, together with the standard packages those lines need.  The retained environments print in this document as Theorem 1 in order of appearance, whereas the article prints the same items with the numbering of its own full document.  The complete native source of the article as distributed in the arXiv e-print for this exact version is bundled unchanged as \texttt{sources/199096/source0.tex}.
\begin{theorem}\label{th1}
Let $\omega \in A_p$ and $y \in W^{2,p}_\omega(0,T)$. The uniform error of the L1 scheme satisfies
\begin{equation}
    |D^\alpha y(t_n) - \delta_\tau^\alpha y(t_n)| \le C \tau^{2-\alpha-1/p} \Lambda(\omega, \tau) \|y\|_{W^{2,p}_\omega},
\end{equation}
where
\begin{equation}
    \Lambda(\omega, \tau) = \sup_{0 \le j \le N-1} \tau^{-1/q} W_j = \sup_{j} \left( \frac{1}{\tau} \int_{t_j}^{t_{j+1}} \omega(s)^{-\frac{1}{p-1}} ds \right)^{\frac{p-1}{p}}.
\end{equation}
\end{theorem}

\begin{proof}
We consider a uniform grid $t_n = n\tau$ for $n=0, \dots, N$, with step size $\tau = T/N$. The L1 scheme is given by replacing $y'(s)$ with the backward difference quotient $u_j$:
\begin{equation} \label{eq:L1}
    \delta_\tau^\alpha y(t_n) = \frac{1}{\Gamma(1-\alpha)} \sum_{j=0}^{n-1} u_j \int_{t_j}^{t_{j+1}} (t_n - s)^{-\alpha} ds,
\end{equation}
where $u_j = \frac{y(t_{j+1}) - y(t_j)}{\tau}$. Let us define the seminorm
\begin{equation}
    [y]_{W^{2,p}_\omega} := \|y''\|_{L^p_\omega}.
\end{equation}
Since $[y]_{W^{2,p}_\omega} \le \|y\|_{W^{2,p}_\omega}$, we will derive the bounds in terms of this seminorm. Let us denote $R_n(y) := |D^\alpha y(t_n) - \delta_\tau^\alpha y(t_n)|$. Subtracting \eqref{eq:L1} from the exact derivative, we have
\begin{equation} \label{eq:error_rep}
    R_n(y) = \frac{1}{\Gamma(1-\alpha)} \left|\sum_{j=0}^{n-1} \int_{t_j}^{t_{j+1}} (t_n - s)^{-\alpha} (y'(s) - u_j) ds \right|.
\end{equation}
Since $y'$ is an absolutely continuous function, using the Taylor series expansion with the integral form of the remainder, for all $s \in (t_j, t_{j+1})$, we have
\begin{equation}
    |y'(s) - u_j| = \left|y'(s) - \frac{1}{\tau}\int_{t_j}^{t_{j+1}} y'(\xi) d\xi\right| = \left|\tau^{-1} \int_{t_j}^{t_{j+1}} \int_{\xi}^s y''(\eta) d\eta d\xi\right| \le \int_{t_j}^{t_{j+1}} |y''(\eta)| d\eta.
\end{equation}
Let $Y_j := \int_{t_j}^{t_{j+1}} |y''(\eta)| d\eta$. Thus, $\sup_{s \in (t_j, t_{j+1})} |y'(s) - u_j| \le Y_j$. Since the only singularity occurs in the term $j = n - 1$ with $n \ge 1$, this term has to be estimated separately. By direct calculation, we have
\begin{align}
    \left|\int_{t_{n-1}}^{t_{n}} (t_n - s)^{-\alpha} (y'(s) - u_{n-1}) ds\right| &\le \int_{t_{n-1}}^{t_n} (t_n - s)^{-\alpha} |y'(s) - u_{n-1}| ds \nonumber \\
    &\le Y_{n-1} \int_{t_{n-1}}^{t_n} (t_n - s)^{-\alpha} ds = \frac{\tau^{1-\alpha}}{1-\alpha} Y_{n-1}.
\end{align}
We estimate the remaining part as follows. Since $\int_{t_j}^{t_{j+1}} (y'(s) - u_j) ds = 0$, and $(t_n - t_{j+1})^{-\alpha}$ are constant, we have
\begin{align}
    \left|\int_{t_j}^{t_{j+1}} (t_n - s)^{-\alpha} (y'(s) - u_j) ds\right| &= \left|\int_{t_j}^{t_{j+1}} \left[ (t_n - s)^{-\alpha} - (t_n - t_{j+1})^{-\alpha} \right] (y'(s) - u_j) ds \right| \nonumber \\
    &\le \int_{t_j}^{t_{j+1}} |(t_n - s)^{-\alpha} - (t_n - t_{j+1})^{-\alpha}| |y'(s) - u_j| ds.
\end{align}
Let $k(s) = (t_n-s)^{-\alpha}$. From the Mean Value Theorem, we have $|k(s) - k(t_{j+1})| = |k'(\xi)||s-t_{j+1}|$ for some $\xi\in[t_j, t_{j+1}]$. Since $k'(s) = \alpha(t_n-s)^{-\alpha-1}$, we have $\max_{\xi\in[t_j, t_{j+1}]}k'(\xi) = \alpha(t_n-t_{j+1})^{-\alpha-1}$ and $|k(s) - k(t_{j+1})| \leq\tau\alpha(t_n-t_{j+1})^{-\alpha-1}$. Hence, we have
\begin{align}
    \left|\int_{t_j}^{t_{j+1}} (t_n - s)^{-\alpha} (y'(s) - u_j) ds \right| &\le \int_{t_j}^{t_{j+1}} |(t_n - s)^{-\alpha} - (t_n - t_{j+1})^{-\alpha}| |y'(s) - u_j| ds \nonumber \\
    &\le \alpha \tau^2 (t_n - t_{j+1})^{-\alpha-1} Y_j.
\end{align}
From H\"older's inequality with conjugates $p, q$ ($1/p + 1/q = 1$), we have
\begin{align}
    Y_j &= \int_{t_j}^{t_{j+1}} |y''(s)| \omega(s)^{1/p} \omega(s)^{-1/p} ds \nonumber \\
    &\le \left( \int_{t_j}^{t_{j+1}} |y''(s)|^p \omega(s) ds \right)^{1/p} \left( \int_{t_j}^{t_{j+1}} \omega(s)^{-q/p} ds \right)^{1/q} \nonumber \\
    &= \|y''\|_{L^p_\omega(t_j, t_{j+1})} \|\omega^{-1/p}\|_{L^q(t_j, t_{j+1})}.
\end{align}
Let us denote $W_j := \|\omega^{-1/p}\|_{L^q(t_j, t_{j+1})}$. Plugging all these calculations into \eqref{eq:error_rep}, we have
\begin{equation}
    R_n(y) \le C \left[ \tau^{1-\alpha} W_{n-1} \|y''\|_{L^p_\omega(t_{n-1}, t_{n})} + \sum_{j=0}^{n-2} \tau^2 (t_n - t_{j+1})^{-\alpha-1} W_j \|y''\|_{L^p_\omega(t_j, t_{j+1})} \right].
\end{equation}
Applying H\"older's inequality to the sum in the above inequality, we have
\begin{equation}
    \sum_{j=0}^{n-2} A_j B_j \le \left( \sum_{j=0}^{n-2} A_j^q \right)^{1/q} \left( \sum_{j=0}^{n-2} B_j^p \right)^{1/p},
\end{equation}
where $A_j =\tau^2 (t_n - t_{j+1})^{-\alpha-1} W_j$ and $B_j = \|y''\|_{L^p_\omega(t_j, t_{j+1})}$. Note that $(\sum B_j^p)^{1/p} = \|y''\|_{L^p_\omega(0,t_n)} \le \|y\|_{W^{2,p}_\omega}$. Now, let us analyze the term:
\begin{equation}
    \left( \sum_{j=0}^{n-2} A_j^q \right)^{1/q} = \left( \sum_{j=0}^{n-2} \left[ \tau^2 (t_n - t_{j+1})^{-\alpha-1} W_j \right]^q \right)^{1/q}.
\end{equation}
To do this, we introduce
\begin{equation}
    \Lambda(\omega, \tau) := \sup_{0 \le j \le N-1} \tau^{-1/q} W_j = \sup_{j} \left( \frac{1}{\tau} \int_{t_j}^{t_{j+1}} \omega(s)^{-\frac{1}{p-1}} ds \right)^{\frac{p-1}{p}}.
\end{equation}
Then $W_j \le \tau^{1/q} \Lambda(\omega, \tau)$. Thus, we have
\begin{align}
    \left( \sum_{j=0}^{n-2} A_j^q \right)^{1/q} &\le \tau^2 \tau^{1/q} \Lambda(\omega, \tau) \left( \sum_{j=0}^{n-2} (t_n - t_{j+1})^{-(\alpha+1)q} \right)^{1/q} \nonumber \\
    &\le \tau^2 \tau^{1/q} \Lambda(\omega, \tau) \left( \frac{1}{\tau} \int_0^{t_n-\tau} (t_n - s)^{-(\alpha+1)q} ds \right)^{1/q} \nonumber \\
    &\le \tau^2 \Lambda(\omega, \tau) \left( \int_0^{t_n-\tau} (t_n - s)^{-(\alpha+1)q} ds \right)^{1/q} \nonumber \\
    &\le C_2 \Lambda(\omega, \tau) \tau^{2-(\alpha+1) + 1/q} = C_2 \Lambda(\omega, \tau) \tau^{2-\alpha-1/p}.
\end{align}
This concludes the proof.
\end{proof}

\end{document}
